% Section 8.5, from lectures/L08/appendix/d_exercises.md. The note under the lecture's intro, which
% exercise is the Kontroll, is folded into the aside, as in Chapter 1.

\section{Övningar}
\label{c8:sec:exercises}\label{c8:app:d}

\begin{aside}
\textbf{Så kontrollerar du ditt arbete.} Övningarna märkta \textbf{Program} och \textbf{Kontroll}
görs i Microchip Studio: skriv programmet, bygg och stega, och läs registren i Processor Status och
portarna i I/O-fönstret. Skriv ned din förutsägelse innan du stegar.

Övriga övningar görs med papper och penna, utan miniräknare. Lösningsförslag finns i
bilaga~\ref{app:answers}.

Varje övning är märkt med sin sort. Det finns en \textbf{Kontroll} i det här kapitlet,
övning~\ref{c8:ex:5}.
\end{aside}

\exercise{Binär addition för hand}{Räkna för hand}
\label{c8:ex:1}
Addera binärt, kolumn för kolumn med minnessiffror, och ange resultatet i en byte, hexadecimalt,
och värdet på \code{C} efteråt.
\begin{parts}
  \item \code{0b00101101 + 0b00010111}
  \item \code{0x9C + 0x7A}
  \item \code{0x0F + 0x01}. Vilken flagga, utöver \code{Z}, \code{N} och \code{C}, sätts här, och
    varför?
\end{parts}

\exercise{Mätaren slår om}{Förståelse}
\label{c8:ex:2}
\code{r16} innehåller 253 och \code{C} är 0. Programmet kör \code{inc r16} fem gånger i rad.
\begin{parts}
  \item Vilket värde har \code{r16} efter varje \code{inc}?
  \item Efter vilken \code{inc} är \code{Z = 1}? Vad är \code{Z} när alla fem är körda?
  \item Vad är \code{C} efter den femte \code{inc}? Varför?
  \item Hur skulle du ändra programmet så att \code{C} blir 1 när räknaren slår om från 255 till 0?
\end{parts}

\exercise{Negativa tal}{Räkna för hand}
\label{c8:ex:3}
\begin{parts}
  \item Skriv $-1$, $-10$, $-100$ och $-128$ som byte i tvåkomplement, hexadecimalt. Visa metoden
    \enquote{invertera och lägg till 1} för $-10$.
  \item Vilka tal med tecken är \code{0xF0}, \code{0x81} och \code{0x7F}?
  \item Vad blir \code{neg r16} om \code{r16} är \code{0x01}, \code{0x80} respektive \code{0x00}?
    Ett av svaren är överraskande. Vilket, och varför?
\end{parts}

\exercise{Addera en konstant}{Program}
\label{c8:ex:4}
\begin{parts}
  \item Skriv ett program som laddar 40 i \code{r16} och sedan adderar 25 till det, utan att
    använda ett andra register.
  \item Lägg till instruktioner som laddar 30 i \code{r17} och drar ifrån 100.
  \item Förutsäg \code{r16}, \code{r17} och \code{C} när programmet är klart. Simulera och
    kontrollera.
\end{parts}

\exercise{Flaggorna}{Kontroll}
\label{c8:ex:5}
\emph{Förutsäg för hand, kontrollera i simulatorn, förklara skillnaderna.}
\begin{parts}
  \item Fyll i tabellen \textbf{för hand}: resultatet hexadecimalt och flaggorna \code{C},
    \code{Z}, \code{N}, \code{V} och \code{S} efter varje operation. Operationerna görs var för
    sig, inte efter varandra.
\end{parts}

\begin{filltable}{@{}p{14em}C{5em}YYYYY@{}}
  \thead{Operation} & \thead{Resultat} & \thead{C} & \thead{Z} & \thead{N} & \thead{V} &
    \thead{S} \\
  \midrule
  \code{0x50 + 0x50} (\code{add}) & & & & & & \fillrule
  \code{0xF0 + 0x20} (\code{add}) & & & & & & \fillrule
  \code{0x30 - 0x30} (\code{sub}) & & & & & & \fillrule
  \code{0x00 - 0x01} (\code{sub}) & & & & & & \fillrule
  \code{0x80 - 0x01} (\code{sub}) & & & & & & \fillrule
  \code{0xFF + 1} med \code{inc}, när \code{C} är 1 före & & & & & & \\
\end{filltable}

\begin{parts}[start=2]
  \item Skriv ett program som gör de sex operationerna i ordning, med en etikett efter varje, som
    i \code{flags_demo.asm} (\secref{c8:a7}). Sätt \code{C} till 1 före den sista med
    instruktionen \code{sec} (\emph{set carry}).
  \item Kör fram till varje etikett och jämför SREG med tabellen. Stämde allt? Förklara varje ruta
    där du hade fel, och säg vilken regel i \secref{c8:app:a} du hade missat.
\end{parts}

\exercise{Masker}{Räkna för hand}
\label{c8:ex:6}
\code{r16} innehåller \code{0b11010011}. Vad blir \code{r16} efter var och en av instruktionerna
nedan? Varje instruktion utgår från \code{0b11010011}.
\begin{parts}
  \item \code{andi r16, 0xF0}
  \item \code{ori r16, 0x0C}
  \item \code{eor r16, r24}, där \code{r24} innehåller \code{0xFF}
  \item \code{com r16}. Jämför med \textbf{c)}.
\end{parts}

\exercise{Välj rätt mask}{Konstruktion}
\label{c8:ex:7}
Skriv en instruktion, eller två om det behövs, som gör följande med \code{r16} och lämnar alla
andra bitar orörda:
\begin{parts}
  \item ettställer bit~0 och bit~1;
  \item nollställer bit~7;
  \item inverterar bit 4--7;
  \item nollställer alla bitar utom bit~3;
  \item nollställer bit 0--3 och ettställer bit~7.
\end{parts}

\exercise{Skift och rotation}{Räkna för hand}
\label{c8:ex:8}
\begin{parts}
  \item \code{r16} innehåller \code{0b00000110}. Vad blir den efter två \code{lsl}? Vilket tal har
    den multiplicerats med?
  \item \code{r17} innehåller \code{0b10000001}. Vad blir \code{r17} och \code{C} efter en
    \code{lsr}?
  \item \code{r18} innehåller \code{0xF8}, det vill säga $-8$. Vad blir den efter en \code{asr},
    och efter en \code{lsr}? Vilket av svaren är $-8 / 2$?
  \item \code{r19} innehåller \code{0b11000000} och \code{C} är 0. Vad blir \code{r19} och
    \code{C} efter var och en av tre \code{rol} i rad?
\end{parts}

\exercise{Lysdioder på port B}{Program}
\label{c8:ex:9}
\begin{parts}
  \item Skriv ett program som gör port B till utport och tänder lysdioderna på PB0, PB2 och PB4
    med en enda \code{out}.
  \item Skriv om programmet så att samma tre lysdioder tänds med \code{sbi}, en i taget, i
    stället.
  \item Lägg till instruktioner som släcker PB2 med \code{cbi}, och sedan tänder PB7 med mönstret
    läs, ändra, skriv (\secref{c8:c5}).
  \item Förutsäg \code{PORTB} efter programmet, binärt. Kontrollera i I/O-fönstret.
\end{parts}

\exercise{Läs programmet}{Förståelse}
\label{c8:ex:10}
Förutsäg \code{PORTB}, binärt, efter var och en av de markerade raderna.

\begin{asmcode}
.def leds = r16
.def temp = r17

main:
    ser temp
    out DDRB, temp
    ldi leds, 0b00001111
    out PORTB, leds                 ; (1)
    lsl leds
    lsl leds
    out PORTB, leds                 ; (2)
    sbi PORTB, 0                    ; (3)
    com leds
    out PORTB, leds                 ; (4)
    cbi PORTB, 7                    ; (5)
end:
    rjmp end
\end{asmcode}

\exercise{Multiplicera med tio}{Program, Fördjupning}
\label{c8:ex:11}
Processorn har ingen instruktion som multiplicerar med 10, men $10x = 8x + 2x$, och multiplikation
med 2 och 8 är skift.
\begin{parts}
  \item Skriv ett program som laddar ett tal $x$ i \code{r16} och räknar ut $10x$ i \code{r16}, med
    \code{mov}, \code{lsl} och \code{add}.
  \item Prova med $x = 7$. Vad blir resultatet, decimalt och hexadecimalt?
  \item Upp till vilket $x$ ger programmet rätt svar? Vad händer för större $x$?
\end{parts}

\exercise{När är det overflow?}{Förståelse, Fördjupning}
\label{c8:ex:12}
\begin{parts}
  \item \code{0x40 + 0x40} ger \code{V = 1}, men \code{0xC0 + 0xC0} ger \code{V = 0}. Förklara båda
    genom att läsa talen med tecken.
  \item Vad blir \code{C} i de två fallen? Förklara genom att läsa talen utan tecken.
  \item Formulera med egna ord en regel för när en addition ger \code{V = 1}.
\end{parts}
